\begingroup
\renewcommand{\arraystretch}{1.45}
\fancyhead[L]{\small ZUR THEORIE DER NEUEN TRANSZENDENTEN. V.}

\Needspace{22\baselineskip}
\par\medskip\Needspace{5\baselineskip}\pdfbookmark[1]{Art. 3}{transscendentes-v-art-3}\centerline{[3.]}\nobreak\smallskip
\[
\begin{gathered}
P(x^3,1)=P,\qquad Q(x^3,1)=Q,\qquad R(x^3,1)=R\\[8pt]
P(x,1)=p,\qquad Q(x,1)=q,\qquad R(x,1)=r
\end{gathered}
\]
\[
\begin{gathered}
pqr^3.P(x^3,y^3)=pqR.P(x,y)^3\\[8pt]\quad{}+3PQR.P(x,y).Q(x,y)^2\\[8pt]
pq^3r.P(x^3,y^3)=prQ.P(x,y)^3\\[8pt]\quad{}-3PQR.P(x,y).R(x,y)^2
\end{gathered}
\]
\[
\begin{gathered}
3pPQR=r^3Q-q^3R\\[8pt]
3qPQR=r^3P-p^3R\\[8pt]
3rPQR=p^3Q-q^3P
\end{gathered}
\]
\[
3PPQQ=PPqq+ppQQ+ppqq
\]
\[
\begin{gathered}
3pP=\frac{r^3}{R}-\frac{q^3}{Q}=\frac{3Q^3}{q}-\frac{3R^3}{r}\\[8pt]
3qQ=\frac{r^3}{R}-\frac{p^3}{P}=\frac{3P^3}{p}-\frac{3R^3}{r}\\[8pt]
3rR=\frac{p^3}{P}-\frac{q^3}{Q}=\frac{3Q^3}{q}-\frac{3P^3}{p}
\end{gathered}
\]

\Needspace{32\baselineskip}
\[
\begin{gathered}
\frac{r^4+3R^4}{rR}=\frac{q^4+3Q^4}{qQ}=\frac{p^4+3P^4}{pP}\\[8pt]
=4+24xx+24x^6+24x^8\\[8pt]\quad{}+48x^{14}+24x^{18}+24x^{24}+48x^{26}\\[8pt]
=8P(xx,1)P(x^6,1)-4Q(xx,1)Q(x^6,1)\\[8pt]\quad{}\ ?\quad\text{würde dann sein}\\[8pt]
=4\surd(pp+qq)(PP+QQ)-4\surd pqPQ\\[8pt]
=2\surd(pp+qq)(PP+QQ)+2\surd(pp-qq)(PP-QQ)
\end{gathered}
\]



\Needspace{13\baselineskip}
\SourceLeaf{481}


\[
\begin{gathered}
\frac{dp}{p}-\frac{dP}{P}=+\tfrac12qrQR.\frac{dx}{x}\\[8pt]
\frac{dq}{q}-\frac{dQ}{Q}=-\tfrac12prPR.\frac{dx}{x}\\[8pt]
\frac{dr}{r}-\frac{dR}{R}=+\tfrac12pqPQ.\frac{dx}{x}
\end{gathered}
\]


\Needspace{29\baselineskip}
\par\medskip\Needspace{5\baselineskip}\pdfbookmark[1]{Art. 4}{transscendentes-v-art-4}\centerline{[4.]}\nobreak\smallskip
Wie
\[
x^{\frac14}(y^{\frac12}+y^{-\frac12})+x^{\frac94}(y^{\frac32}+y^{-\frac32})+.\ +\ .\ .\ .
\]
durch $R(x,y)$ bezeichnet ist, so wollen wir noch
\[
x^{\frac14}(y^{\frac12}-y^{-\frac12})-x^{\frac94}(y^{\frac32}-y^{-\frac32})+.\ -\ .\ .\ .
\]
durch $S(x,y)$ bezeichnen. Dann gilt
\[
iS(x,y)=R(x,-y)
\]

\Needspace{19\baselineskip}
\[
\begin{gathered}
S(x,y)=\frac{\surd[Q(x,1)^2.P(x,y)^2-P(x,1)^2.Q(x,y)^2]}{R(x,1)}\\[8pt]
=\frac{\surd[P(x,1)^2.R(x,y)^2-R(x,1)^2.P(x,y)^2]}{Q(x,1)}
\end{gathered}
\]
\[
\begin{gathered}
R(x,y).S(x,y)=Q(xx,1).S(xx,yy)\\[8pt]
R(x,y)^2+S(x,y)^2=2P(xx,1).R(xx,yy)\\[8pt]
R(x,y)^2-S(x,y)^2=2R(xx,1).P(xx,yy)
\end{gathered}
\]

\Needspace{32\baselineskip}
\[
\begin{gathered}
R(x,y).R(x,z)+S(x,y).S(x,z)\\[8pt]\quad{}=2R(xx,yz).P\Bigl(xx,\frac yz\Bigr)\\[8pt]
R(x,y).R(x,z)-S(x,y).S(x,z)\\[8pt]\quad{}=2R\Bigl(xx,\frac yz\Bigr).P(xx,yz)\\[8pt]
R(x,y).S(x,z)+S(x,y).R(x,z)\\[8pt]\quad{}=2S(xx,yz).Q\Bigl(xx,\frac yz\Bigr)\\[8pt]
R(x,y).S(x,z)-S(x,y).R(x,z)\\[8pt]\quad{}=2S\Bigl(xx,\frac zy\Bigr).Q(xx,yz)
\end{gathered}
\]

\Needspace{22\baselineskip}
Setzt man
\[
\Sigma e^{-\theta(k+\alpha)^2}=(\theta,\alpha)
\]
wo $k$ alle ganzen Zahlen bedeutet, so sind die Sätze enthalten in
\[
\begin{gathered}(\theta,2\alpha).(\theta,2\beta)
\\[8pt]\quad{}=(2\theta,\alpha+\beta).(2\theta,\alpha-\beta)
\\[8pt]\quad{}+(2\theta,\alpha+\beta+\tfrac12).(2\theta,\alpha-\beta+\tfrac12)\end{gathered}
\]



\providecommand{\GaussCircleArrow}{\ensuremath{\mathord{\vcenter{\hbox{\setlength{\unitlength}{1pt}\begin{picture}(10,11)\linethickness{.55pt}\put(3.3,3.3){\circle{5.2}}\put(5.1,5.1){\line(1,1){3.2}}\put(5.6,8.3){\line(1,0){2.7}}\put(8.3,5.6){\line(0,1){2.7}}\end{picture}}}}}}

\Needspace{8\baselineskip}
\SourceLeaf{482}
auch ist
\[
(\theta,\alpha)=(\theta,-\alpha)
\]
\clearpage
\endgroup
